\section{A Gentle Map Of The Destination}

Blaze is a protocol for proving evaluations of large committed tensors while
reading only a small number of symbols.  The final construction uses several
pieces at once: finite-field arithmetic, multilinear extensions, commitments,
linear error-correcting codes, sumcheck, row compression, interleaving, RAA
encoding, and BaseFold-style proximity checks.

This paper builds those pieces bottom-up.  First we separate the proof models
and oracle interfaces, so words like ``oracle,'' ``implicit input,'' and
``commitment'' have precise meanings before Blaze starts composing them.  Then
we define the algebraic objects: fields, tensors, multilinear extensions,
commitments, codes, sumcheck, and folding.  After that we explain Blaze's
code-independent compression step: encode many row slices, combine their
committed row words into an openable folded tensor $c_\star$, and bind
$c_\star$ back to the committed matrix.  Only after that do we instantiate the
linear code with RAA and explain the correlated checks that make the proof
sound.

The one high-level idea to keep in mind is this: the protocol does not run code
checks and tensor-evaluation checks in isolation.  The opened positions,
folded terminal values, and multilinear claims must all refer to the same
committed objects.  Much of the backend exists to maintain that correlation.
